% ============================================================
% 中文学位论文模板（thesis-cn）
% ------------------------------------------------------------
% 基于 ctexbook 文档类，默认使用 TeX Live 自带的 fandol 字体，
% 无需在系统中安装额外字体。编译引擎：xelatex。
% 目录结构：
%   main.tex            主文件（封面 / 声明 / 摘要 / 目录 / 全局设置）
%   chapters/chap01.tex 第一章：绪论
%   chapters/chap02.tex 第二章：方法与实现（示例）
%   chapters/conclusion.tex 结论与展望
% 修改指南：
%   1. 在"封面信息"处填写学校、题目、姓名、学号、指导教师等；
%   2. 摘要、Abstract 与关键词直接替换示例文字；
%   3. 各章在 chapters/ 目录中编辑，新增章节时
%      复制一份 chap 文件并在 main.tex 中 \include；
%   4. 参考文献在文末 thebibliography 中按引用顺序补充；
%   5. 学校若有专用封面 / 声明格式，以其官方 Word/PDF 为准替换本框架。
% ============================================================
\documentclass[
  UTF8,
  fontset=fandol,        % TeX Live 自带开源字体，跨平台可编译
  zihao=-4,              % 正文字号：小四
  a4paper,
  oneside,               % 单面排版（双面打印可改 twoside）
]{ctexbook}

% ---- 常用宏包 ----
\usepackage{amsmath, amssymb}             % 数学公式
\usepackage{graphicx}                     % 插图
\usepackage{booktabs}                     % 三线表
\usepackage{multirow}                     % 表格合并单元格
\usepackage{caption}                      % 图表标题样式
\captionsetup{font=small, labelfont=bf}
\usepackage{geometry}
\geometry{left=3cm, right=2.6cm, top=3cm, bottom=2.8cm}
\usepackage{setspace}
\setstretch{1.5}                          % 正文 1.5 倍行距（常见要求）
\usepackage[
  colorlinks=true,
  linkcolor=black,
  citecolor=black,
  urlcolor=black,
]{hyperref}                               % 打印版链接全部黑色

% 数学环境符号约定：向量用粗斜体
\newcommand{\vect}[1]{\boldsymbol{#1}}

% ---- 章节标题样式：一级标题黑体居中（常见本科论文要求）----
\ctexset{
  chapter = {
    format = \centering\heiti\zihao{3},
    nameformat = {},
    titleformat = {},
    number = \arabic{chapter},          % 第 1 章（可用 \chinese{chapter} 改为中文数字）
    beforeskip = 18pt,
    afterskip = 24pt,
  },
  section = {
    format = \heiti\zihao{4},
  },
  subsection = {
    format = \heiti\zihao{-4},
  },
}

% 图表标题改为"图 1-1 / 表 2-1"式编号（章号-图号）
\renewcommand{\thefigure}{\thechapter-\arabic{figure}}
\renewcommand{\thetable}{\thechapter-\arabic{table}}
\renewcommand{\theequation}{\thechapter-\arabic{equation}}

% ============================================================
% 封面信息：在这里填写你的论文信息（全文唯一需要集中修改的地方）
% ============================================================
\newcommand{\SchoolName}{示例大学}
\newcommand{\ThesisTitle}{基于示例技术的示例系统设计与实现}
\newcommand{\ThesisAuthor}{张~三}
\newcommand{\StudentNumber}{2022010000}
\newcommand{\Major}{计算机科学与技术}
\newcommand{\Advisor}{李~四~教授}
\newcommand{\ThesisDate}{2026年6月}

\begin{document}

% ==================== 封面 ====================
\begin{titlepage}
  \begin{center}
    {\zihao{1}\heiti \SchoolName}\\[0.8em]
    {\zihao{2}\heiti 本科毕业论文（设计）}\\[4em]
    {\zihao{2}\heiti \ThesisTitle}\\[6em]
    \zihao{3}
    \renewcommand{\arraystretch}{1.6}
    \begin{tabular}{rl}
      {\heiti 姓~~~~~~名：} & \ThesisAuthor \\
      {\heiti 学~~~~~~号：} & \StudentNumber \\
      {\heiti 专~~~~~~业：} & \Major \\
      {\heiti 指导教师：} & \Advisor \\
    \end{tabular}\\[6em]
    {\zihao{3} \ThesisDate}
  \end{center}
\end{titlepage}

% ==================== 原创性声明与授权页 ====================
% 声明文字为通用示例，正式提交前请替换为学校官方文本。
\chapter*{原创性声明}
\addcontentsline{toc}{chapter}{原创性声明}
本人郑重声明：所呈交的论文是本人在指导教师指导下独立进行研究
所取得的成果。除了文中特别加以标注引用的内容外，本论文不包含
任何其他个人或集体已经发表或撰写的成果作品。对本文的研究做出
重要贡献的个人和集体，均已在文中以明确方式标明。
本人完全意识到本声明的法律后果由本人承担。

\vspace{4em}
\begin{flushright}
  论文作者签名：\underline{\hspace{6em}} \quad 日期：\underline{\hspace{6em}}
\end{flushright}

\chapter*{版权使用授权书}
\addcontentsline{toc}{chapter}{版权使用授权书}
本人完全了解学校有关保留、使用学位论文的规定，同意学校保留并
向国家有关部门或机构送交论文的复印件和电子版，允许论文被查阅
和借阅。本人授权学校可以将本学位论文的全部或部分内容编入有关
数据库进行检索，可以采用影印、缩印或扫描等复制手段保存和汇编
本学位论文。

\vspace{4em}
\begin{flushright}
  \begin{tabular}{l}
    论文作者签名：\underline{\hspace{6em}} \quad 日期：\underline{\hspace{6em}} \\[1.2em]
    指导教师签名：\underline{\hspace{6em}} \quad 日期：\underline{\hspace{6em}} \\
  \end{tabular}
\end{flushright}

% ==================== 中文摘要 ====================
\chapter*{摘\quad 要}
\addcontentsline{toc}{chapter}{摘\quad 要}
% 摘要：目的—方法—结果—结论，一般 300～500 字
在此撰写中文摘要。简要说明研究背景与问题、所提出的方法、
主要实验结果以及结论与创新点。摘要应独立成文，
不出现图表、公式编号与参考文献引用。

\vspace{1em}
\noindent\textbf{关键词：}关键词一；关键词二；关键词三；关键词四

% ==================== 英文摘要 ====================
\chapter*{Abstract}
\addcontentsline{toc}{chapter}{Abstract}
Write the English abstract here, mirroring the content of the Chinese
abstract. Keep it self-contained: state the problem, the approach,
the main results, and the conclusion.

\vspace{1em}
\noindent\textbf{Key words:} Keyword One; Keyword Two; Keyword Three; Keyword Four

% ==================== 目录 ====================
\tableofcontents

% ==================== 正文 ====================
% 第一章：绪论
% 论文第一章一般交代研究背景、现状与本文工作。
% 本章演示了常用的章节层级与图表引用方式。
\chapter{绪\quad 论}

\section{研究背景与意义}
% 用 1～2 段说明课题来源、要解决的实际问题与研究意义
在此阐述研究背景与课题意义。可以先介绍问题的实际背景，
再指出现有方法存在的不足，从而引出本文的研究动机。

\section{国内外研究现状}
\subsection{方向一的研究现状}
按技术路线或时间顺序综述代表性工作，引用示例：统计学习理论
框架见文献~\cite{vapnik2000}，深度学习方法的综述见
文献~\cite{lecun2015}，Transformer 架构由文献~\cite{vaswani2017} 提出。

\subsection{方向二的研究现状}
继续综述与本课题相关的其他方向，并总结现有工作的不足。

\section{主要研究内容}
% 用一段话 + 条目列出本文的三到四项工作
本文的主要研究内容包括：
\begin{itemize}
  \item 研究内容一：对问题进行建模与形式化描述；
  \item 研究内容二：设计并实现示例系统；
  \item 研究内容三：通过实验验证方法的有效性。
\end{itemize}

\section{论文组织结构}
% 逐章一句话说明全文结构
本文共分为三章：第一章介绍研究背景与主要内容；
第二章介绍方法设计与实验分析；第三章总结全文并展望未来工作。

% 第二章：方法与实现（示例章）
% 演示公式、表格、插图的规范用法，写作时替换为实际内容。
\chapter{方法设计与实验分析}

\section{总体设计}
% 给出系统/方法的总体框架，配一张框架图
介绍总体设计方案。核心目标函数如式~\eqref{eq:objective}：
\begin{equation}
  \min_{\vect{\theta}} \;
  \mathcal{L}(\vect{\theta}) =
  \frac{1}{N} \sum_{i=1}^{N}
  \ell \bigl( f(\vect{x}_i; \vect{\theta}),\, y_i \bigr)
  + \lambda \, \lVert \vect{\theta} \rVert_2^2
  \label{eq:objective}
\end{equation}
其中 $\ell(\cdot)$ 为损失函数，$\lambda$ 为正则化系数。

\section{实验设置}
% 数据集、评价指标、对比方法、实现细节（硬件/超参数）
说明数据集划分、评价指标与实现细节。

\section{结果与分析}
% 主实验表格：三线表是学位论文的标准表格样式
表~\ref{tab:compare} 给出了各方法在基准数据集上的对比结果。

\begin{table}[!htbp]
  \centering
  \caption{各方法在基准数据集上的性能对比}
  \label{tab:compare}
  \begin{tabular}{lccc}
    \toprule
    方法 & 准确率（\%） & $F_1$（\%） & 参数量（M） \\
    \midrule
    基线方法 A & 85.2 & 84.1 & 21.3 \\
    基线方法 B & 87.6 & 86.9 & 34.8 \\
    本文方法   & \textbf{90.4} & \textbf{89.7} & 28.5 \\
    \bottomrule
  \end{tabular}
\end{table}

% 插图示例：把图片文件放到本目录（支持 PNG/PDF/JPG），替换文件名即可
% \begin{figure}[!htbp]
%   \centering
%   \includegraphics[width=0.8\linewidth]{framework.png}
%   \caption{系统总体框架图}
%   \label{fig:framework}
% \end{figure}

对实验结果进行分析：先描述现象，再解释原因，
最后回应第一章提出的研究问题。

% 第三章：结论与展望
% 总结全文工作、指出局限性、展望未来方向，一般 1～2 页。
\chapter{总结与展望}

\section{工作总结}
% 逐条对应绪论中的研究内容，概述做了什么、得到了什么结论
本文围绕示例问题开展了研究，主要工作总结如下：
\begin{itemize}
  \item 对问题进行了形式化建模，明确了优化目标；
  \item 设计并实现了示例系统，验证了方案的可行性；
  \item 实验结果表明，本文方法在主要指标上优于基线方法。
\end{itemize}

\section{不足与展望}
% 坦诚指出工作的局限性，并提出后续可改进的方向
受时间与条件限制，本文工作仍存在不足：样本规模有限、
对某些边界情况的讨论不够充分。未来可以从以下方向继续改进：
扩大实验规模、引入更多约束条件、探索方法的可解释性。


% ==================== 参考文献 ====================
% 按 GB/T 7714 风格手工排版，条目按正文引用顺序编号；
% 文献较多时建议改用 BibTeX（\bibliographystyle{...} + \bibliography{refs}）。
% 注意：book 类下 thebibliography 自带“参考文献”章标题，
% 这里不要再手写 \chapter*{参考文献}，否则会出现重复标题。
\begin{thebibliography}{9}
\addcontentsline{toc}{chapter}{参考文献}
\setstretch{1.2}
\bibitem{vapnik2000}
VAPNIK V N. 统计学习理论[M]. 张学工, 译. 北京: 清华大学出版社, 2000.

\bibitem{lecun2015}
LECUN Y, BENGIO Y, HINTON G. Deep learning[J]. Nature, 2015,
521(7553): 436-444.

\bibitem{vaswani2017}
VASWANI A, SHAZEER N, PARMAR N, et al. Attention is all you
need[C]//Advances in Neural Information Processing Systems. 2017: 5998-6008.
\end{thebibliography}

% ==================== 致谢 ====================
\chapter*{致\quad 谢}
\addcontentsline{toc}{chapter}{致\quad 谢}
在此撰写致谢。感谢指导教师在选题、研究与写作过程中的悉心指导，
感谢同组同学与家人的支持与帮助。致谢内容应真诚得体，
避免过度夸张。

\end{document}
